\documentclass[11pt]{article}
\usepackage[T1]{fontenc}
\usepackage[utf8]{inputenc}
\usepackage{lmodern,amsmath,amssymb,microtype,xcolor}
\usepackage[margin=.9in]{geometry}
\providecommand{\tightlist}{\setlength{\itemsep}{0pt}\setlength{\parskip}{0pt}}
\usepackage[colorlinks=true,urlcolor=blue!60!black]{hyperref}
\setlength{\parindent}{0pt}
\setlength{\parskip}{.5em}
\setlength{\emergencystretch}{2em}
\author{Wentao Li\\\small Prepared with assistance from OpenAI Codex}
\date{November 22, 2026}
\title{Two prescribed targets for rearranging relative density}
\hypersetup{pdfauthor={Wentao Li},pdftitle={Two prescribed targets for rearranging relative density}}
\begin{document}
\maketitle
Brech, Brendle and Telles asked in Question 56 whether a target set of fewer than continuum many interior densities forces the relative-density permutation number to equal the continuum. Valderrama asked in Question 4.5 whether the number can be smaller than the continuum when every interior target is allowed. We construct a ccc extension in which two prescribed targets already suffice. \href{https://arxiv.org/abs/2410.21102}{Brech, Brendle and Telles, 2024, Definition 44 and Question 56} · \href{https://doi.org/10.1112/jlms.70300}{2025 journal version} · \href{https://arxiv.org/abs/2506.21059}{Valderrama, 2025, Definition 1.6 and Question 4.5} · \href{https://doi.org/10.1007/s00153-026-01041-3}{2026 journal version}

\textbf{Theorem.} Fix distinct \(u,v\in(0,1)\) in a ground model with \(\mathfrak c\ge\aleph_2\). There is a ccc extension satisfying \[\mathfrak{dd}^{\rm rel}_{(0,1),\{u,v\}}=
\mathfrak{dd}^{\rm rel}_{(0,1),(0,1)}=\aleph_1<\mathfrak c.\] For \(A\subseteq B\subseteq\omega\) with \(B\) infinite, \(d_B(A)\) is the limit of \(|A\cap N|/|B\cap N|\), when it exists. A witnessing family of permutations must send each pair of original density \(p\in(0,1)\) to a pair of existing relative density different from \(p\) in the stated target set. The theorem includes the fixed targets \(1/3,1/2\).

\textbf{Random conditional probabilities.} Work first over one ground model. Index a compact cube \(K=[0,1]^D\) by all its pairs \(D=\{(A,B): A\subseteq B\subseteq\omega,\ B\text{ infinite}\}\). Fix one target \(t\in(0,1)\). For each integer \(N\ge2\), give the \(N\) points below \(N\) independent levels \(L\) and independent symmetric binary marks. Use \(\Pr(L=j)=2^{-j-1}\) for \(j\ge0\). A point's weight is \(N^{2L}\) times \(t\) or \(1-t\) according to the mark. Define \(q_N(A,B)\) as the ratio of total weight on \(A\cap N\) to total weight on \(B\cap N\), using zero when the denominator vanishes. These random vectors have Borel probability laws on \(K\) (the finite sample has countably many possibilities). Take a weak-star cluster law \(\mu\) along a nonprincipal ultrafilter on \(N\), using compactness of the Radon probability measures on \(K\). Equivalently integrate each continuous function, take the ultralimit and use the Riesz representation theorem.

The law has two properties. First, almost surely on each fixed finite list of coordinates, the vector can be approximated arbitrarily well by conditional probabilities of a real probability distribution assigning strictly positive probabilities to the infinite atoms of the finite Boolean algebra of the sets in the list and zero to its finite atoms. To see this, delete the union \(F\) of those finite atoms. For a fixed denominator \(B\) let \(b=|B\cap N|\to\infty\). With probability tending to one the maximum level on \(B\setminus F\) strictly exceeds every level in \(F\). For a proof, first bound all levels on \(F\) by a large fixed integer except on an event of arbitrarily small probability, then let \(b\) tend to infinity. When the strict gap holds, the relative mass of \(F\) in that denominator is at most \(|F|\max(t,1-t)/(N^2\min(t,1-t))\). This works for all finitely many denominators together. After deletion every infinite atom is represented for all large \(N\). The distance to the closed set of feasible limiting vectors therefore goes to zero in probability, and this closed set has full measure under the limiting finite-dimensional law.

Second, if \(d_B(A)=p\in(0,1)\), then \(\mu\{q:q(A,B)=t\}>0\). Set \(j=\lfloor\log_2 b\rfloor\). Require exactly one point of \(A\cap N\) at level \(j\) with weight multiplier \(t\), exactly one point of \((B\setminus A)\cap N\) there with multiplier \(1-t\), and all other \(B\) points below \(j\). With \(a=|A\cap N|\), the event has probability \[a(b-a)\,(2^{-j-2})^2(1-2^{-j})^{b-2}.\] Its lower limit is strictly positive since \(a/b\to p\) and \(1\le b2^{-j}<2\). On that event the top two weights determine the ratio up to \(O(1/N)\) since at most \(N\) lower weights occur. For every closed neighborhood of \(t\) the limiting coordinate law gives at least the same positive lower bound by weak convergence along the ultrafilter. Shrinking these neighborhoods gives the asserted atom.

Force with the probability algebra for countably many independent copies of this law for each of the two targets. Each random coordinate is a name for a real, obtained for instance from its rational sublevel sets. For every fixed ground pair of interior density, with Boolean value one some copy for each target hits that target exactly. For each copy every old finite list is feasible in the sense above with Boolean value one. Both assertions hold separately for every member of the corresponding ground set, which suffices for the forcing universal quantifier. This use of the probability algebra does not require a single point in the old cube satisfying uncountably many requirements.

\textbf{Realizing one feasible vector by an actual enumeration.} In this intermediate extension let \(q\) be one vector with the finite feasibility property. A condition consists of a finite injective list \(s\) of naturals and a finite set \(F\) of triples \((A,B,e)\) formed from old pairs and rational tolerances \(0<e\le1\). For each triple let \(m_B\) be the number of entries of \(s\) in \(B\). Require \(m_B>0\), the empirical ratio to differ from \(q(A,B)\) by at most \(e/2\), and \[\sum_{(A,B,e)\in F}\frac{32}{e^2m_B}<1/8.\] The empty list with no tests is allowed. Extensions append to the list, enlarge \(F\), and keep the deviation at most \(e\) at \textbf{every} intermediate prefix for each old test. This is a partial order by transitivity.

Here is the extension calculation. It applies even to the union of two test sets for the same initial list, where the displayed sum is less than \(1/4\). Include any finite list of new tests as well. Choose a distribution on the infinite Boolean atoms that approximates all the target conditional ratios within \(e/8\) for their respective tolerances. Simulate independent draws of these atom labels. Realize each draw by a distinct, previously unused integer of the specified atom. Each denominator is hit infinitely often almost surely. Along hits of \(B\) the marked successes for \(A\) are independent Bernoulli variables of parameter \(p'\) within \(e/8\) of \(q(A,B)\).

For their centered partial sums \(S_l\), the probability that \(|S_l|>(e/2)(m_B+l)\) at some \(l\) is at most \(32/(e^2m_B)\). Indeed use Kolmogorov's maximal inequality first up to \(m_B\) and then up to \(2^{h+1}m_B\) for \(2^hm_B\le l\le2^{h+1}m_B\), and sum the geometric bounds, using variance at most one per draw. Outside this event the initial error at most \(em_B/2\), the drift at most \(el/8\) and the centered error together stay below \(e(m_B+l)\) (for \(l=0\) the initial condition itself suffices). A union bound gives positive probability of meeting every old prefix test throughout the simulation. The strong law also holds for every listed pair. Eventually the final deviations of all old and new tests are less than \(e/3\), and their counts are large enough for any requested budget. Stop at such a prefix. If a specified integer must appear, reserve it in choosing representatives and append it last unless it is already present. By taking the counts sufficiently large this last change preserves the prefix tests and leaves all final deviations at most \(e/2\). Thus adding any test and including any integer are dense requirements.

Conditions with the same list are linked by this calculation. The forcing is therefore \(\sigma\)-linked and ccc. The generic list enumerates \(\omega\) bijectively. The inverse of the enumerating map is a permutation \(\pi\) with \(d_{\pi[B]}(\pi[A])=q(A,B)\) for all old pairs: add the test at successively smaller rational tolerances, then use the intermediate-prefix requirement. Use a countable finite-support iteration of these ccc forcings to realize all copies. Every old pair of interior density now reaches each of the two targets under a member of a countable family. At least one differs from its initial density.

\textbf{Iteration and lower bound.} Repeat this whole ccc step in a finite-support iteration of length \(\omega_1\). Each pair of subsets of \(\omega\) in the final extension has a name supported in a bounded set of stages, by ccc and the regularity of \(\omega_1\). Existence and value of a density are absolute for the same pair between stages. The union of the stage families thus gives the two upper bounds. Ground \(\aleph_2\) and its distinct reals survive. For the lower bounds, no countable family of permutations can change the density of every fair-coin subset of \(\omega\) relative to \(B=\omega\): by the strong law one common set and all its countably many images have density \(1/2\). This proves both equalities.

\textbf{Scope and priority.} Brech, Brendle and Telles's Theorem 55 gives continuum for the analogous series problem with a small set of finite target sums. Valderrama's prior forcing gives a small family for endpoint targets. The interior-target argument uses ties at the highest random level to obtain exact atoms, followed by a forcing that realizes conditional probabilities even on sparse denominators. Provisional significance \textbf{70/100} if correct and new. The general rearrangement-number comparison in Question 51 remains separate. Independent proof and priority review are pending.
\end{document}
